% GENERATED FILE. Edit canonical lessons, then run npm run generate:books.
% canonical-source-hash: fnv1a64:bc2daa26ff6b3fb8
% canonical-lessons: TE-C454-pailatu, TE-C454-antarjatiya, TE-C454-svadesam, TE-C454-ravana, TE-C454-saganampu

\chapter{Flying Abroad}
\label{ch:te-flying-abroad}
\hlchaptermodality{454}

\begin{chapteropening}
\textbf{By the end of this chapter:} \emph{I can say a pilot, international, one's homeland, transport and to see someone off.}

Complete the last lesson of chapter 454: I can say a pilot, international, one's homeland, transport and to see someone off.
\end{chapteropening}

\section[\texorpdfstring{\te{పైలటు} (pailaṭu)}{పైలటు (pailaṭu)}]{\te{పైలటు} (pailaṭu) — a pilot}

\label{lesson:TE-C454-pailatu}



\begin{quote}
Before the new one: say the Telugu for to float away, then the Telugu for to be swept away.
\end{quote}

\subsection*{You'll want to know: \te{పైలటు}}
\textbf{\te{పైలటు}} — \emph{pailaṭu} — \textquotedblleft{}a pilot\textquotedblright{}.

\begin{quote}
\te{విమానం} \te{దిగే} \te{ముందు} \te{పైలటు} \te{మాతో} \te{మాట్లాడాడు}.
\end{quote}

\begin{quote}
\emph{vimānaṁ digē mundu pailaṭu mātō māṭlāḍāḍu.}
\end{quote}

\begin{quote}
Before the plane landed, the pilot spoke to us.
\end{quote}

Like other English loans it takes -\te{ఉ}: \te{పైలటు}. A woman pilot is also \te{పైలటు}.

\begin{cousinweb}
\textbf{\te{పైలటు}} (\emph{pailaṭu}) is a loan from English \emph{pilot}, with the Telugu final -\te{ఉ}.
\end{cousinweb}

\begin{grammarlens}[title={What you've built}]
One more word for getting around.
\end{grammarlens}

\subsection*{Your turn}
\begin{itemize}
\raggedright
  \item \emph{Say it:} \emph{pailaṭu}
  \item \emph{Say it:} \emph{pailaṭu}, once more
  \item \emph{Say it:} \emph{vimānaṁ digē mundu pailaṭu mātō māṭlāḍāḍu.} — Before the plane landed, the pilot spoke to us.
  \item \emph{Recall:} say the Telugu for to float away, then the Telugu for to be swept away, then say \emph{pailaṭu} again
\end{itemize}

\begin{tcolorbox}[breakable,colback=teal!4,colframe=teal!35!black,title={Before you move on}]
What does \te{పైలటు} mean? (\textquotedblleft{}A pilot\textquotedblright{}.) Say it once more.
\end{tcolorbox}
\section[\texorpdfstring{\te{అంతర్జాతీయ} (antarjātīya)}{అంతర్జాతీయ (antarjātīya)}]{\te{అంతర్జాతీయ} (antarjātīya) — international}

\label{lesson:TE-C454-antarjatiya}



\begin{quote}
Before the new one: say the Telugu for to be swept away, then the Telugu for a pilot.
\end{quote}

\subsection*{You'll want to know: \te{అంతర్జాతీయ}}
\textbf{\te{అంతర్జాతీయ}} — \emph{antarjātīya} — \textquotedblleft{}international\textquotedblright{}.

\begin{quote}
\te{మా} \te{నగరంలో} \te{కొత్త} \te{అంతర్జాతీయ} \te{విమానాశ్రయం} \te{ఉంది}.
\end{quote}

\begin{quote}
\emph{mā nagaraṁlō kotta antarjātīya vimānāśrayaṁ undi.}
\end{quote}

\begin{quote}
There is a new international airport in our city.
\end{quote}

An adjective: \te{అంతర్జాతీయ} \te{విమానం}, an international flight. It stands before its noun and does not change.

\begin{cousinweb}
\textbf{\te{అంతర్జాతీయ}} (\emph{antarjātīya}) is a learned word built from Sanskrit: \emph{antar}, \textquotedblleft{}between\textquotedblright{}, \emph{jāti}, \textquotedblleft{}birth, a people, a nation\textquotedblright{}, and the adjective ending \emph{-īya} (Monier-Williams): between the nations.
\end{cousinweb}

\begin{grammarlens}[title={What you've built}]
One more word for getting around.
\end{grammarlens}

\subsection*{Your turn}
\begin{itemize}
\raggedright
  \item \emph{Say it:} \emph{antarjātīya}
  \item \emph{Say it:} \emph{antarjātīya}, once more
  \item \emph{Say it:} \emph{mā nagaraṁlō kotta antarjātīya vimānāśrayaṁ undi.} — There is a new international airport in our city.
  \item \emph{Recall:} say the Telugu for to be swept away, then the Telugu for a pilot, then say \emph{antarjātīya} again
\end{itemize}

\begin{tcolorbox}[breakable,colback=teal!4,colframe=teal!35!black,title={Before you move on}]
What does \te{అంతర్జాతీయ} mean? (\textquotedblleft{}International\textquotedblright{}.) Say it once more.
\end{tcolorbox}
\section[\texorpdfstring{\te{స్వదేశం} (svadēśaṁ)}{స్వదేశం (svadēśaṁ)}]{\te{స్వదేశం} (svadēśaṁ) — one's homeland}

\label{lesson:TE-C454-svadesam}



\begin{quote}
Before the new one: say the Telugu for a pilot, then the Telugu for international.
\end{quote}

\subsection*{You'll want to know: \te{స్వదేశం}}
\textbf{\te{స్వదేశం}} — \emph{svadēśaṁ} — \textquotedblleft{}one's homeland\textquotedblright{}.

\begin{quote}
\te{చాలా} \te{సంవత్సరాల} \te{తరువాత} \te{అతను} \te{స్వదేశానికి} \te{తిరిగి} \te{వచ్చాడు}.
\end{quote}

\begin{quote}
\emph{cālā saṁvatsarāla taruvāta atanu svadēśāniki tirigi vaccāḍu.}
\end{quote}

\begin{quote}
After many years he came back to his homeland.
\end{quote}

\te{స్వ}, own, + \te{దేశం}, country. Its opposite is \te{విదేశం}, abroad.

\begin{cousinweb}
\textbf{\te{స్వదేశం}} (\emph{svadēśaṁ}) is borrowed from Sanskrit \emph{svadeśa}, \textquotedblleft{}one's own place or country\textquotedblright{} (Monier-Williams), from \emph{sva}, \textquotedblleft{}own\textquotedblright{}, and \emph{deśa}, \textquotedblleft{}a place, a country\textquotedblright{}.
\end{cousinweb}

\begin{grammarlens}[title={What you've built}]
One more word for getting around.
\end{grammarlens}

\subsection*{Your turn}
\begin{itemize}
\raggedright
  \item \emph{Say it:} \emph{svadēśaṁ}
  \item \emph{Say it:} \emph{svadēśaṁ}, once more
  \item \emph{Say it:} \emph{cālā saṁvatsarāla taruvāta atanu svadēśāniki tirigi vaccāḍu.} — After many years he came back to his homeland.
  \item \emph{Recall:} say the Telugu for a pilot, then the Telugu for international, then say \emph{svadēśaṁ} again
\end{itemize}

\begin{tcolorbox}[breakable,colback=teal!4,colframe=teal!35!black,title={Before you move on}]
What does \te{స్వదేశం} mean? (\textquotedblleft{}One's homeland\textquotedblright{}.) Say it once more.
\end{tcolorbox}
\section[\texorpdfstring{\te{రవాణా} (ravāṇā)}{రవాణా (ravāṇā)}]{\te{రవాణా} (ravāṇā) — transport}

\label{lesson:TE-C454-ravana}



\begin{quote}
Before the new one: say the Telugu for international, then the Telugu for one's homeland.
\end{quote}

\subsection*{You'll want to know: \te{రవాణా}}
\textbf{\te{రవాణా}} — \emph{ravāṇā} — \textquotedblleft{}transport\textquotedblright{}.

\begin{quote}
\te{నగరంలో} \te{ప్రజా} \te{రవాణా} \te{బాగుంది}.
\end{quote}

\begin{quote}
\emph{nagaraṁlō prajā ravāṇā bāgundi.}
\end{quote}

\begin{quote}
Public transport in the city is good.
\end{quote}

\te{ప్రజా} \te{రవాణా}, public transport; \te{సరుకు} \te{రవాణా}, the carrying of goods.

\begin{cousinweb}
\textbf{\te{రవాణా}} (\emph{ravāṇā}) is usually traced to Persian \emph{ravāna}, \textquotedblleft{}setting out, dispatched\textquotedblright{}, which came into Telugu through Hindustani; sending goods on their way became transport.
\end{cousinweb}

\begin{grammarlens}[title={What you've built}]
One more word for getting around.
\end{grammarlens}

\subsection*{Your turn}
\begin{itemize}
\raggedright
  \item \emph{Say it:} \emph{ravāṇā}
  \item \emph{Say it:} \emph{ravāṇā}, once more
  \item \emph{Say it:} \emph{nagaraṁlō prajā ravāṇā bāgundi.} — Public transport in the city is good.
  \item \emph{Recall:} say the Telugu for international, then the Telugu for one's homeland, then say \emph{ravāṇā} again
\end{itemize}

\begin{tcolorbox}[breakable,colback=teal!4,colframe=teal!35!black,title={Before you move on}]
What does \te{రవాణా} mean? (\textquotedblleft{}Transport\textquotedblright{}.) Say it once more.
\end{tcolorbox}
\section[\texorpdfstring{\te{సాగనంపు} (sāganampu)}{సాగనంపు (sāganampu)}]{\te{సాగనంపు} (sāganampu) — to see off}

\label{lesson:TE-C454-saganampu}



\begin{quote}
Before the new one: say the Telugu for one's homeland, then the Telugu for transport.
\end{quote}

\subsection*{You'll want to know: \te{సాగనంపు}}
\textbf{\te{సాగనంపు}} — \emph{sāganampu} — \textquotedblleft{}to see off\textquotedblright{}.

\begin{quote}
\te{మేము} \te{అత్తను} \te{విమానాశ్రయంలో} \te{సాగనంపాం}.
\end{quote}

\begin{quote}
\emph{mēmu attanu vimānāśrayaṁlō sāganampāṁ.}
\end{quote}

\begin{quote}
We saw our aunt off at the airport.
\end{quote}

Past (I …): \te{సాగనంపాను}, \emph{sāganampānu}.

\te{సాగను}, to go on, + \te{అంపు}, to send: to send someone on their way. Telugu families often go to the station to see a guest off.

\begin{cousinweb}
\textbf{\te{సాగనంపు}} (\emph{sāganampu}) is native Dravidian: \te{సాగను}, from \te{సాగు}, \textquotedblleft{}to go on, to proceed\textquotedblright{} (DEDR 2430), and \te{అంపు}, \textquotedblleft{}to send\textquotedblright{} (DEDR 329).
\end{cousinweb}

\begin{grammarlens}[title={What you've built}]
One more word for getting around.
\end{grammarlens}

\subsection*{Your turn}
\begin{itemize}
\raggedright
  \item \emph{Say it:} \emph{sāganampu}
  \item \emph{Say it:} \emph{sāganampu}, once more
  \item \emph{Say it:} \emph{mēmu attanu vimānāśrayaṁlō sāganampāṁ.} — We saw our aunt off at the airport.
  \item \emph{Recall:} say the Telugu for one's homeland, then the Telugu for transport, then say \emph{sāganampu} again
\end{itemize}

\begin{tcolorbox}[breakable,colback=teal!4,colframe=teal!35!black,title={Before you move on}]
What does \te{సాగనంపు} mean? (\textquotedblleft{}To see off\textquotedblright{}.) Say it once more.
\end{tcolorbox}
